% TT75040C
% Stromunfälle
% 14.0
% 2013-11-30
{
.. _m-stromunfall:

}
{
\begin{itemize}
  -   **Eigenschutz!** Bevor der Patienten berührt wird,
  muss man sich vergewissern, dass der Patient
  keinen Kontakt mehr mit dem Stromkreis hat! \EE{Strom
  abschalten (lassen)!} (Vgl. \FullPrettyRef{topic:gefahrenbereich-strom})
  \begin{itemize}
    -   Gefahrenbreich beachten
    -   Sicherheitsabstand
    -   Schrittspannung und Spannungskegel
    beachten 
    -   Bahnanlagen sperren/sichern lassen
    -   Schalter ausschalten, (Haupt-)Sicherung
    auslösen, Starkstrom/Hochspannung durch Fachmann  abschalten bzw.
    erden lassen
    -   Strom gegen Wiedereinschalten sichern!
    -   Patient mit isolierendem Material aus dem Stromkreis befreien
  \end{itemize}
  -   Beurteilung der vitalen Bedrohung.
  \par |TxBeiVit|
  |TxMassVitMK|
  -   **Monitoring und
  Reanimationsbereitschaft**: auf evtl.
  Herzrhythmusstörungen vorbereitet sein
  (bis zu 24 h Latenz)
  -   Genaue Patientenuntersuchung ${\rightarrow}$
  Eintritts- und Austrittsmarke
  -   Hospitalisierung auch bei symptomlosen
  Patienten zur Überwachung
  -   Stromverbrennungen wie Verbrennungen versorgen
\end{itemize}
}
